\documentclass{article}

\title{HW 11 Solutions}
\author{Principles of Microeconomics}
\date{\today}

\begin{document}

\maketitle

\begin{enumerate}

\item

	\begin{enumerate}
	
	\item
	
		\begin{itemize}
		
		\item Police protection is a club good because it’s excludable but not rival. You can deny police protection to certain people, but one person’s use of police protection doesn’t usually diminish another person’s. 
		
		\item Snow plowing is a public good because it’s not excludable or rival. You can’t prevent people from using a plowed street, and one person’s use of the plowed street doesn’t usually diminish another person’s.
		
		\item Education is a private good because it’s excludable and rival. You can prevent people who don’t pay for an education from receiving it, and since classrooms can only fit so many kids, one person’s consumption of education prevents another person from consuming it.
		
		\item Rural roads are public goods because they’re not excludable or rival. You can’t prevent people from using a rural road, and there’s no traffic in rural areas, so one person’s use doesn’t diminish another person’s
		
		\item City streets are common resources because they’re rival but not excludable. They’re not excludable because you can’t prevent people from using them, but they’re rival because there’s so much traffic in cities that one person’s use diminishes another person’s by adding to congestion.
		
		\end{itemize}
	
	\item If the private sector doesn’t produce enough of a good because it has a positive externality, the government can increase production by producing the good itself. This is the case with education and police protection. Sometimes the public uses a government provided public good so much that it becomes a common resource. This is the case with city streets.
	
	\end{enumerate}
	
\item 

	\begin{enumerate}
	
	\item Public goods typically have positive externalities, so the free-market quantity is less than the efficient quantity. For example, national defense provides protection for all of the people inside a country — not just the people who pay for it. As a result, it has a positive externality.
	
	\item Common resources typically have negative externalities, so the free-market quantity is greater than the efficient quantity. For example, when someone uses a city street during rush hour, they create a negative externality because they create congestion that slows other people down.
	
	\end{enumerate}
	
\item

	\begin{enumerate}
	
	\item Fredo is a free rider.
	
	\item The government can provide non-excludable goods and pay for them by increasing taxes.
	
	\item The private market can provide TV in an excludable way like through a cable, satellite, or streaming service. That way people have to pay for using the service. Any of these solutions turn TV from a public good into a club good.
	
	\end{enumerate}
	
\item

	\begin{enumerate}
	
	\item It’s a public good because it’s not excludable or rival. It’s free, so everyone can use it without paying, and it doesn’t have many users so no one person’s use slows down another person’s use.
	
	\item It’s a common resource because it’s become rival. One person’s use diminishes another’s.
	
	\item People might overuse the internet because they don’t account for the costs they create for others when they decide how much to use it. One way to correct that is by adding a paywall and thereby turning it into a private good.
	
	\end{enumerate}

\end{enumerate}

\end{document}